\documentclass[../../../include/open-logic-section]{subfiles}

\begin{document}

\iftag{FOL}
      {\olfileid{fol}{seq}{rul}}
      {\olfileid{pl}{seq}{rul}}

\olsection{规则与\usetoken{P}{derivation}}

在下文中，令 $\Gamma, \Delta, \Pi, \Lambda$ 代表有限的语句序列。
% Part: first-order-logic
% Chapter: sequent-calculus
% Section: rules-and-proofs

\begin{defn}[相继式]
\emph{相继式}是形如
\[
\Gamma \Sequent \Delta
\]
的表达式，其中 $\Gamma$ 和 $\Delta$ 是语言 $\Lang L$ 中的有限语句序列（可能为空）。$\Gamma$ 称为\emph{前件}，而 $\Delta$ 称为\emph{后件}。
\end{defn}

\begin{explain}
相继式的直观含义是：如果前件中的所有语句都成立，那么后件中至少有一个语句成立。也就是说，如果 $\Gamma = \tuple{!A_1, \dots, !A_m}$ 且
$\Delta = \tuple{!B_1, \dots, !B_n}$，那么 $\Gamma \Sequent \Delta$ 成立当且仅当
\[
(!A_1 \land \cdots \land !A_m) \lif (!B_1 \lor \cdots \lor
!B_n)
\]
成立。有两种特殊情况：一种是 $\Gamma$ 为空，另一种是 $\Delta$ 为空。当 $\Gamma$ 为空，即 $m = 0$ 时，$\quad \Sequent \Delta$ 成立当且仅当 $!B_1 \lor \dots \lor !B_n$ 成立。当 $\Delta$ 为空，即 $n = 0$ 时，$\Gamma \Sequent \quad$ 成立当且仅当 $\lnot(!A_1 \land \dots \land !A_m)$ 成立。我们称一个相继式是\emph{有效的}，当且仅当其对应的语句是有效的。
\end{explain}

如果 $\Gamma$ 是一个语句序列，我们用 $\Gamma, !A$ 表示将 $!A$ 附加到 $\Gamma$ 右端的结果（而用 $!A, \Gamma$ 表示将 $!A$ 附加到 $\Gamma$ 左端）。如果 $\Delta$ 也是一个语句序列，那么 $\Gamma, \Delta$ 表示两个序列的拼接。

\begin{defn}[初始相继式]
\emph{初始相继式}是
\iftag{prvFalse,prvTrue}{以下形式之一的相继式：
  \begin{enumerate}
    \item $!A \Sequent !A$
    \tagitem{prvTrue}{$\quad \Sequent \ltrue$}{}
    \tagitem{prvFalse}{$\lfalse \Sequent \quad$}{}
  \end{enumerate}}
{形如 $!A \Sequent !A$ 的相继式}，其中 $!A$ 是语言中的任意语句。
\end{defn}

在相继式演算中，推导是由相继式构成的特定树形结构，其中最顶层的相继式是初始相继式，并且如果一个相继式排在一个或两个相继式下方，它必须由某条推理规则正确地得出。$\Log{LK}$ 的规则分为两大类型：\emph{逻辑}规则和\emph{结构}规则。逻辑规则得名于下层相继式中包含 $!A$ 和/或 $!B$ 的语句的\emph{主联结词}。每条规则都有两个版本，一个用于推导包含该逻辑算子的语句位于左方的相继式，另一个则用于推导该语句位于右方的相继式。

\end{document}